\Subhead{3.1.2}{Статистическое и временное усреднение. Эргодичность. Измерение вероятностей и средних значений}

Ансамблевое среднее $\langle X(t)\rangle$ сравнивает много опытов в одно время; временное $\overline X_T=T^{-1}\int_0^TX(t)\dd t$ использует одну запись. Эргодичность по среднему означает сходимость второго к первому в оговорённом смысле. Это дополнительное свойство, а не следствие слова «стационарный».

Для стационарного процесса
\begin{equation}
 \operatorname{Var}\overline X_T=\frac{2}{T}\int_0^T\left(1-\frac\tau T\right)C_X(\tau)\dd\tau.
\end{equation}
Формула получается из двойного интеграла ковариации по квадрату $[0,T]^2$. Если $\int_0^\infty|C_X(\tau)|\dd\tau<\infty$, дисперсия среднего стремится к нулю. При ненулевом интеграле корреляции и $T$ много больше времени корреляции $\tau_c=\int_0^\infty C_X(\tau)\dd\tau/C_X(0)$ имеем $\operatorname{Var}\overline X_T\simeq2\sigma^2\tau_c/T$.

Контрпример: $X(t)=A$, где $A$ — случайная константа опыта. Процесс стационарен, но временное среднее равно случайному $A$, а не $\langle A\rangle$. Для измерения вероятности события считают долю времени/опытов, когда оно выполнено. У независимых испытаний дисперсия частоты $p(1-p)/N$; корреляции меняют точность оценки через сумму ковариаций. Положительные длительные корреляции обычно уменьшают эффективное число независимых отсчётов; антикорреляции могут его увеличить. Нельзя оценивать погрешность длинной шумовой записи, просто считая все цифровые отсчёты независимыми.
